\section{Introduction}

Several distinct questions are often grouped under the heading of quantum communication through gravity. One may ask how quantum matter radiates into a quantized gravitational field, how propagating gravitons correlate separated matter systems, how a quantum gravitational-wave mode is absorbed by a detector, or whether an effective gravitational interaction channel can preserve entanglement. Each ingredient now has substantial prior literature.

Closed or conserved branch-dependent sources and their emitted coherent graviton states have been analyzed directly, including the relation between the coherent-state distance of two radiation histories and vacuum which-path decoherence \cite{Matsui2026}. Classical conserved stress tensors driving the quantized gravitational field produce coherent states whose expectation value reproduces the retarded classical waveform \cite{LagaSuyama2026}. Propagating gravitons have been used to generate delayed matter entanglement \cite{TrengganaZen2026}, while resonant matter systems have been treated as receivers of nonclassical gravitational-wave states with open-system noise included \cite{ToccaceloEtAl2026,TobarEtAl2024}. Gravitational quantum communication has also been phrased directly in terms of entanglement-breaking channels and thermal thresholds \cite{MariEtAl2026,ToccaceloAndersenBrask2025,MikiLiChen2026}. Quantum photon--graviton conversion, cavity-QED transduction, and stimulated graviton absorption/emission have likewise been proposed \cite{IkedaEtAl2025,AraniLamineSoda2026,Schutzhold2025}.

Our question is narrower. We ask what happens when the interfaces usually idealized away are kept explicit in one source-to-readout normalization:
\begin{equation}
\boxed{
\begin{aligned}
&\text{local equal-charge encoding}
\to
\text{conserved radiator}
\to
\text{gravitational output port}\\
&\hspace{1.2cm}\to
\text{retarded free-space mode}
\to
\text{receiver memory}
\to
\text{readout}.
\end{aligned}}
\label{eq:chain}
\end{equation}
The objective is not to introduce a new information-theory criterion. The entanglement-breaking structure of one-mode Gaussian channels is established \cite{Holevo2008EB}, as is the sensitivity of non-Gaussian entanglement to attenuation and amplification \cite{FilippovZiman2014}. The contribution sought here is instead a source-resolved gravitational link calculation in which the physical meaning and normalization of each efficiency and noise term remain visible. Figure~\ref{fig:link-budget} summarizes the architecture and separates the coherent link backbone from the later memory-noise and readout tests.

\begin{figure}[t]
\centering
\resizebox{\textwidth}{!}{%
\begin{tikzpicture}[
  >=Latex,
  node distance=1.25cm,
  box/.style={draw, rounded corners=2pt, align=center, minimum height=1.55cm, minimum width=3.1cm, inner sep=6pt},
  arr/.style={->, line width=0.8pt},
  lab/.style={align=center, font=\small}
]

\node[box] (source) {\textbf{Equal-charge local code}\\
$V_{\mathcal C}^\dagger Q_A V_{\mathcal C}=q_A I$\\
conserved source; handoff $T_*$};

\node[box, right=of source] (gmode) {\textbf{Normalized graviton mode}\\
complete encoder + tail waveform\\
$\int |f(t)|^2dt=1$};

\node[box, right=of gmode] (rport) {\textbf{Receiver gravitational port}\\
retarded wave-zone access\\
$R/c$ causal lower bound};

\node[box, right=of rport] (memory) {\textbf{Receiver memory}\\
$\tau_c=\eta_Q^{\rm link}\mathcal T_f$\\
$\Delta_{\rm mem}=\tau_c-m_c$};

\node[box, right=of memory] (readout) {\textbf{Accessible register}\\
readout $(\tau_r,m_r)$\\
$\Delta_{\rm acc}=\tau_r\Delta_{\rm mem}-m_r$};

\draw[arr] (source) -- node[lab, above] {$\displaystyle \beta_{g,A}=\frac{\kappa_{g,A}}{\kappa_A}$\\source branching} (gmode);

\draw[arr] (gmode) -- node[lab, above] {$\displaystyle \eta_{\rm store}=\frac{25\mathcal O}{16(kR)^2}$\\mode capture} (rport);

\draw[arr] (rport) -- node[lab, above] {$\displaystyle \beta_{g,B}\mathcal T_f$\\branching $\times$ loading} (memory);

\draw[arr] (memory) -- node[lab, above] {noisy readout} (readout);

\node[lab, below=0.55cm of gmode] (backbone) {$\displaystyle \eta_Q^{\rm link}=\beta_{g,A}\eta_{\rm store}\beta_{g,B}$};

\end{tikzpicture}%
}
\caption{Source-resolved gravitational quantum-link architecture after the local code is instantiated. The three multiplicative factors $\beta_{g,A}$, $\eta_{\rm store}$, and $\beta_{g,B}$ form the coherent link backbone; the fixed waveform contributes $\mathcal T_f\le1$. Noise determines the memory quantum excess, and a separate readout stage determines whether that excess remains accessible. Encoder precursor radiation is retained in the complete waveform $f(t)$ rather than reset at the handoff.}
\label{fig:link-budget}
\end{figure}


A second issue is locality. Gauge-invariant matter operators in gravity require nonlocal gravitational dressing, so the ordinary local-QFT picture of exactly commuting compact subsystems cannot simply be imported into this problem \cite{DonnellyGiddings2016}. Perturbative gravitational splitting nevertheless permits information inside a region to be hidden from exterior measurements, at first order, except through total Poincar\'e charges \cite{DonnellyGiddings2018}. We therefore encode the source qubit into two branches for which the matrix of every total Poincar\'e generator within the logical code is proportional to the identity to the working order, and use a common first-order asymptotic dressing. The operational signal is the later retarded difference field, not an assumed instantaneous branch-dependent Coulomb dressing.

The resulting paper has five logical layers. First, we construct the equal-charge conserved source. Second, we identify the time at which its locally generated bosonic code becomes a fixed channel input. Third, we derive a four-factor coherent link budget. Fourth, we include source and receiver noise and quantify actual entanglement transfer. Fifth, we append a separate noisy readout stage so that memory capture is not confused with accessible quantum reception.
